% !TeX root = ../slides.tex
% =============================================================================
% Bloco de Resultados — seis slides (pontos 1 a 7 de INSTRUCOES.md, seção 3.1)
% Valores: codigo/resultados/*.csv, iguais aos de escrita/06_capitulo/tabelas/.
%
% Estrutura de cada slide: gráfico à esquerda, números em destaque à direita,
% nota com as condições de leitura e, no rodapé, a frase de fechamento.
% =============================================================================

% -----------------------------------------------------------------------------
% R1 — Energia (pergunta A). Fonte: magnitude.csv e estruturada.csv, esquema de
% poda não uniforme (escopo "global"); tabelas energia.tex e custo.tex.
% -----------------------------------------------------------------------------
\begin{frame}
\framesubtitle{Resultados · Pergunta A: custo e energia}
\frametitle{Só a poda estruturada reduziu o consumo de energia}
\begin{columns}[T,onlytextwidth]
\begin{column}{0.58\textwidth}
\begin{tikzpicture}
\begin{axis}[grafico,
  width=6.3cm, height=3.6cm,
  xmin=0, xmax=90, ymin=0, ymax=65,
  xtick={0,10,...,90},
  xticklabel={\pgfmathprintnumber{\tick}\%},
  ytick={0,20,40,60},
  yticklabel={\pgfmathprintnumber{\tick}\%},
  xlabel={taxa de poda},
  legend style={at={(0.03,0.98)}, anchor=north west},
]
  % limiar de detecção do protocolo
  \draw[tintaMuda, dashed, line width=0.6pt] (axis cs:0,10) -- (axis cs:90,10);
  \node[rotuloSec, anchor=south east] at (axis cs:90,10.5) {limiar de detecção (10\%)};

  \addplot[serieDestaque] coordinates {
    (0,0) (10,2.5) (20,5.5) (30,13.5) (40,21.4) (50,28.6) (60,36.2) (70,44.2) (80,51.1) (90,58.7)};
  \addlegendentry{poda estruturada}
  \addplot[serieContextoA] coordinates {
    (0,0) (10,1.7) (20,2.6) (30,2.8) (40,3.7) (50,4.1) (60,5.3) (70,6.2) (80,7.7) (90,9.0)};
  \addlegendentry{poda por magnitude}

  \node[rotulo, anchor=west, font=\fontsize{8}{9}\selectfont\bfseries] at (axis cs:91.5,58.7) {58,7\%};
  \node[rotulo, anchor=west] at (axis cs:91.5,7.6) {9,0\%};
\end{axis}
\end{tikzpicture}
\end{column}
\begin{column}{0.38\textwidth}
\vspace{1mm}
{\fontsize{8.5}{11}\selectfont
\textbf{Com 90\% de taxa de poda}\par\vspace{1.5mm}
\begin{tabular}{@{}l r r@{}}
 & \color{tintaSec}Magnitude & \textbf{Estruturada} \\
\midrule
Parâmetros & \menos 61\% & \textbf{\menos 61\%} \\
Memória    & \color{tintaMuda}0\% & \textbf{\menos 61\%} \\
FLOPs      & \color{tintaMuda}0\% & \textbf{\menos 62\%} \\
Energia    & \color{tintaMuda}n.d. & \textbf{\menos 59\%} \\
Tempo      & \color{tintaMuda}n.d. & \textbf{\menos 56\%} \\
\bottomrule
\end{tabular}}
\vspace{1mm}
\nota{n.d.: não detectável, abaixo do limiar de 10\%.\\
Na poda por magnitude, contam-se os parâmetros não nulos.}
\end{column}
\end{columns}
\vspace{1mm}
\nota{Redução de energia ($\Delta E$) em relação ao modelo base medido na mesma sessão.}
\vfill
\mensagem{A poda por magnitude zera pesos, mas mantém as dimensões dos tensores; a
estruturada remove cabeças de atenção e neurônios do MLP.}
\end{frame}

% -----------------------------------------------------------------------------
% R2 — Perfil de sensibilidade. Fonte: sensibilidade_perfil.csv (taxa-sonda de
% 50%) e sensibilidade_perfil_sonda03.csv (correlação entre sondas); tabela perfil.tex.
% -----------------------------------------------------------------------------
\begin{frame}
\framesubtitle{Resultados · Perfil de sensibilidade por camada}
\frametitle{A camada 0 concentra a sensibilidade do modelo}
\vspace{1.5mm}
\begin{columns}[T,onlytextwidth]
\begin{column}{0.58\textwidth}
\begin{tikzpicture}
\begin{groupplot}[grafico,
  group style={group size=2 by 1, horizontal sep=10mm},
  width=3.2cm, height=2.85cm,
  ybar, /pgf/bar width=4pt,
  xmin=-0.8, xmax=11.8,
  xtick={0,1,...,11},
  x tick style={draw=none},
  ymin=0,
  xlabel={camada},
  title style={font=\fontsize{8.5}{10.5}\selectfont\bfseries, text=tinta,
               anchor=base west, at={(0,1.0)}, inner sep=0pt, yshift=2.5mm},
]
  \nextgroupplot[title={Poda por magnitude},
    ymax=0.45, ytick={0,0.1,0.2,0.3,0.4},
    yticklabel={\pgfmathprintnumber[fixed,precision=1]{\tick}}]
    \addplot[draw=none, fill=contextoB, bar shift=0pt] coordinates {
      (1,0.035) (2,0.093) (3,0.043) (4,0.202) (5,0.130) (6,0.074)
      (7,0.049) (8,0.055) (9,0.061) (10,0.054) (11,0.054)};
    \addplot[draw=none, fill=destaque, bar shift=0pt] coordinates {(0,0.390)};
    \node[rotulo, anchor=south west, font=\fontsize{7.5}{9}\selectfont\bfseries]
      at (axis cs:0.5,0.37) {0,390};

  \nextgroupplot[title={Poda estruturada},
    ymax=4.5, ytick={0,1,2,3,4},
    yticklabel={\pgfmathprintnumber{\tick}}]
    \addplot[draw=none, fill=contextoB, bar shift=0pt] coordinates {
      (1,0.094) (2,0.173) (3,0.183) (4,0.218) (5,0.169) (6,0.138)
      (7,0.118) (8,0.122) (9,0.093) (10,0.138) (11,0.663)};
    \addplot[draw=none, fill=destaque, bar shift=0pt] coordinates {(0,4.009)};
    \node[rotulo, anchor=south west, font=\fontsize{7.5}{9}\selectfont\bfseries]
      at (axis cs:0.5,3.8) {4,01};
    \node[rotulo, anchor=south] at (axis cs:11,0.70) {0,66};
    \node[rotuloSec, anchor=west, align=left] at (axis cs:1.4,2.7)
      {\ppl{} do modelo:\\25,17 $\rightarrow$ 1.386,43};
\end{groupplot}
\end{tikzpicture}
\end{column}
\begin{column}{0.38\textwidth}
\vspace{1mm}
\numero{77\%}{da degradação de podar as 12 camadas é atingida ao podar só a
camada~0 (poda estruturada a 50\%).}
\vspace{3mm}
\numero{42,9\vezes}{é a razão entre a degradação da camada mais sensível e a da
mais tolerante (poda estruturada); 11,3\vezes{} na poda por magnitude.}
\end{column}
\end{columns}
\vspace{1mm}
\nota{$\Delta_\ell$: aumento do logaritmo da \ppl{} ao podar só a camada $\ell$ a
50\%, sem \ft. Escalas verticais diferentes.}
\vfill
\mensagem{A ordenação das camadas se manteve ao repetir o perfil com taxa-sonda de
30\%: correlação de postos de 0,944 (magnitude) e 0,713 (estruturada).}
\end{frame}

% -----------------------------------------------------------------------------
% R3 — Taxas por camada. Fonte: sensibilidade.csv, coluna taxas_por_camada,
% poda estruturada, taxa de poda de 50%, beta = 1,0; tabela alocacao_50.tex.
% -----------------------------------------------------------------------------
\begin{frame}
\framesubtitle{Resultados · Do perfil ao esquema de poda}
\frametitle{O perfil define quanto podar em cada camada}
\begin{columns}[T,onlytextwidth]
\begin{column}{0.58\textwidth}
\begin{tikzpicture}
\begin{axis}[grafico,
  width=7.2cm, height=3.5cm,
  ybar, bar width=9pt,
  xmin=-0.7, xmax=11.7,
  xtick={0,1,...,11},
  x tick style={draw=none},
  ymin=0, ymax=100,
  ytick={0,25,50,75,100},
  yticklabel={\pgfmathprintnumber{\tick}\%},
  xlabel={camada},
]
  \addplot[draw=none, fill=destaque, bar shift=0pt] coordinates {
    (0,1.97) (1,83.64) (2,45.77) (3,43.13) (4,36.17) (5,46.81)
    (6,57.10) (7,67.16) (8,64.49) (9,84.56) (10,57.31) (11,11.90)};
  % esquema de poda uniforme: 50% em todas as camadas
  \draw[contextoA, dashed, line width=0.8pt] (axis cs:-0.7,50) -- (axis cs:11.7,50);
  \node[rotuloSec, anchor=south] at (axis cs:3.5,50.5) {uniforme: 50\%};

  \node[rotulo, anchor=south, font=\fontsize{7.5}{9}\selectfont\bfseries] at (axis cs:0,1.97) {2\%};
  \node[rotulo, anchor=south] at (axis cs:1,83.64) {84\%};
  \node[rotulo, anchor=south] at (axis cs:9,84.56) {85\%};
  \node[rotulo, anchor=south, font=\fontsize{7.5}{9}\selectfont\bfseries] at (axis cs:11,11.90) {12\%};
\end{axis}
\end{tikzpicture}
\end{column}
\begin{column}{0.38\textwidth}
\setbeamerfont{itemize/enumerate body}{size=\fontsize{9.5}{12}\selectfont}
\begin{itemize}
  \setlength{\itemsep}{1.5mm}
  \item Cada camada recebe uma taxa proporcional à sua \textbf{tolerância}: o
        inverso da degradação medida no perfil.
  \item As camadas mais sensíveis são preservadas: \textbf{2\%} na camada~0 e
        \textbf{12\%} na camada~11.
  \item A taxa de poda do modelo não muda: a média das 12 camadas continua
        \textbf{50\%}.
\end{itemize}
\end{column}
\end{columns}
\vspace{1mm}
\nota{Taxa de poda por camada na poda estruturada, para uma taxa de poda de 50\%, com
parâmetro de intensidade \bet{1,0}.}
\vfill
\mensagem{O esquema de poda uniforme aplicaria 50\% também à camada 0, a mais sensível
do modelo.}
\end{frame}

% -----------------------------------------------------------------------------
% R4 — Resultado principal (pergunta B), poda estruturada, 10% a 50%.
% Fonte: estruturada.csv e sensibilidade.csv; tabelas qualidade_estruturada.tex,
% tradeoff.tex (médias) e energia.tex.
% Faixa de 10% a 50%: a mesma dos resumos da Metodologia. Não cortar em 40%.
% -----------------------------------------------------------------------------
\begin{frame}
\framesubtitle{Resultados · Pergunta B: qualidade na poda estruturada}
\frametitle{Análise de sensibilidade: \ppl{} até 4 vezes menor}
\begin{columns}[T,onlytextwidth]
\begin{column}{0.58\textwidth}
\begin{tikzpicture}
\begin{axis}[grafico,
  width=7.0cm, height=3.6cm,
  ymode=log,
  xmin=4, xmax=53, ymin=80, ymax=14000,
  xtick={10,20,30,40,50},
  xticklabel={\pgfmathprintnumber{\tick}\%},
  ytick={100,300,1000,3000,10000},
  yticklabels={100,300,1.000,3.000,10.000},
  xlabel={taxa de poda},
  legend style={at={(0.99,0.03)}, anchor=south east},
]
  \addplot[serieContextoB] coordinates {
    (10,481.13) (20,1484.89) (30,5467.56) (40,10486.71) (50,4573.00)};
  \addlegendentry{sem análise: uniforme}
  \addplot[serieContextoA] coordinates {
    (10,356.20) (20,2102.06) (30,2413.52) (40,2434.92) (50,2402.58)};
  \addlegendentry{sem análise: não uniforme}
  \addplot[serieDestaque] coordinates {
    (10,124.29) (20,371.07) (30,724.53) (40,1301.15) (50,2327.59)};
  \addlegendentry{com análise de sensibilidade}

  % fator entre o melhor esquema sem análise e o esquema com análise
  \draw[tinta, line width=0.5pt] ([xshift=-2.2mm]axis cs:10,356.20) -- ([xshift=-2.2mm]axis cs:10,124.29)
    node[rotulo, midway, left] {2,9\vezes};
  \draw[tinta, line width=0.5pt] ([xshift=2.2mm]axis cs:20,1484.89) -- ([xshift=2.2mm]axis cs:20,371.07)
    node[rotulo, midway, right, font=\fontsize{8.5}{10}\selectfont\bfseries] {4,0\vezes};
  \draw[tinta, line width=0.5pt] ([xshift=2.2mm]axis cs:30,2413.52) -- ([xshift=2.2mm]axis cs:30,724.53)
    node[rotulo, midway, right] {3,3\vezes};
\end{axis}
\end{tikzpicture}
\end{column}
\begin{column}{0.38\textwidth}
\vspace{1mm}
\numero{4,0\vezes}{menor \ppl{} a 20\% de taxa de poda: de 1.484,89 para 371,07.}
\vspace{2.2mm}
\numero{2,5\vezes}{na média geométrica de 10\% a 50\%: de 1.602,62 para 632,46.}
\vspace{2.2mm}
\numerotexto{Mesma energia}{$\Delta E$ médio de 18,0\%, contra 14,3\% e 18,7\%
sem a análise.}
\end{column}
\end{columns}
\vspace{1mm}
\nota{Escala logarítmica, sem \ft{} (modelo base: 25,17). Fatores em relação ao melhor
esquema sem a análise em cada taxa.}
\vfill
\mensagem{Controle: com \bet{0}, o esquema reproduz a poda uniforme até a última
casa decimal. O ganho vem do perfil de sensibilidade.}
\end{frame}

% -----------------------------------------------------------------------------
% R5 — Poda por magnitude. Fonte: magnitude.csv e sensibilidade.csv; tabelas
% qualidade_magnitude.tex e tradeoff.tex (médias).
% A curva azul é beta = 0,5. A contagem "7 de 9" usa o menor valor entre os dois
% betas em cada taxa (beta = 0,5 de 20% a 50%; beta = 1,0 de 60% a 80%).
% -----------------------------------------------------------------------------
\begin{frame}
\framesubtitle{Resultados · Pergunta B: qualidade na poda por magnitude}
\frametitle{O ganho se repetiu na poda por magnitude}
\begin{columns}[T,onlytextwidth]
\begin{column}{0.58\textwidth}
\begin{tikzpicture}
\begin{axis}[grafico,
  width=7.0cm, height=3.6cm,
  ymode=log,
  xmin=4, xmax=53, ymin=20, ymax=1600,
  xtick={10,20,30,40,50},
  xticklabel={\pgfmathprintnumber{\tick}\%},
  ytick={25,50,100,250,500,1000},
  yticklabels={25,50,100,250,500,1.000},
  xlabel={taxa de poda},
  legend style={at={(0.02,0.98)}, anchor=north west},
]
  % modelo base
  \draw[tintaMuda, dashed, line width=0.6pt] (axis cs:4,25.17) -- (axis cs:53,25.17);
  \node[rotuloSec, anchor=south east] at (axis cs:53,25.6) {modelo base: 25,17};

  % diferença a 40% entre o melhor esquema sem análise (uniforme) e o com análise
  \draw[tinta, line width=0.5pt] (axis cs:40,104.82) -- (axis cs:40,57.30);

  \addplot[serieContextoA] coordinates {
    (10,25.40) (20,26.87) (30,40.23) (40,165.20) (50,1096.50)};
  \addlegendentry{sem análise: não uniforme}
  \addplot[serieContextoB] coordinates {
    (10,25.47) (20,27.22) (30,36.42) (40,104.82) (50,473.28)};
  \addlegendentry{sem análise: uniforme}
  \addplot[serieDestaque] coordinates {
    (10,25.44) (20,26.57) (30,31.20) (40,57.30) (50,447.57)};
  \addlegendentry{com análise de sensibilidade}

  \node[rotulo, anchor=north west, font=\fontsize{8.5}{10}\selectfont\bfseries] at (axis cs:40.6,56) {57,30};
\end{axis}
\end{tikzpicture}
\end{column}
\begin{column}{0.38\textwidth}
\numero{7 de 9}{taxas de poda com a menor \ppl{}: todas de 20\% a 80\%.}
\vspace{1.4mm}
\numero{57,30}{contra 104,82 do melhor esquema sem a análise, a 40\%: quase a
metade.}
\vspace{1.4mm}
\numero{\menos 15\%}{na média geométrica de 10\% a 50\%: de 66,00 para 55,80.}
\end{column}
\end{columns}
\vspace{1mm}
\nota{Escala logarítmica, sem \ft. Curva com a análise: \bet{0,5}; a contagem de
7 em 9 considera também \bet{1,0}.}
\vfill
\mensagem{A análise de sensibilidade reduziu a perda de qualidade nas duas estratégias
de poda.}
\end{frame}

% -----------------------------------------------------------------------------
% R6 — Fine-tuning de recuperação, poda estruturada a 30%. Fonte: finetune.csv;
% tabelas finetune.tex e tradeoff.tex (ΔE de 19,1% e 17,7%).
% Referência: modelo base retreinado (18,68), nunca 25,17.
% -----------------------------------------------------------------------------
\begin{frame}
\framesubtitle{Resultados · \textit{Fine-tuning} de recuperação}
\frametitle{Após uma época de \ft, a vantagem se manteve a 30\%}
\begin{columns}[T,onlytextwidth]
\begin{column}{0.58\textwidth}
\begin{tikzpicture}
\begin{axis}[grafico,
  width=3.9cm, height=3.9cm,
  ymode=log,
  xmin=-0.06, xmax=1.06, ymin=14, ymax=3600,
  xtick={0,1},
  xticklabels={sem retreinamento, após 1 época},
  x tick style={draw=none},
  ytick=\empty, ymajorgrids=false,
]
  \addplot[serieContextoA] coordinates {(0,2413.52) (1,104.78)};
  \addplot[serieDestaque] coordinates {(0,724.53) (1,56.08)};
  \addplot[tintaMuda, dashed, line width=0.8pt, mark=*, mark size=1.6pt,
           mark options={solid, fill=tintaMuda, draw=white, line width=0.6pt}]
    coordinates {(0,25.17) (1,18.68)};

  \node[rotulo, anchor=east] at (axis cs:-0.06,2413.52) {2.413,52};
  \node[rotulo, anchor=east] at (axis cs:-0.06,724.53) {724,53};
  \node[rotuloSec, anchor=east] at (axis cs:-0.06,25.17) {25,17};
  \node[rotulo, anchor=west] at (axis cs:1.06,104.78)
    {104,78\hspace{1.5mm}{\color{tintaSec}sem a análise}};
  \node[rotulo, anchor=west] at (axis cs:1.06,56.08)
    {\textbf{56,08}\hspace{1.5mm}{\color{tintaSec}com a análise}};
  \node[rotuloSec, anchor=west] at (axis cs:1.06,18.68)
    {18,68\hspace{1.5mm}modelo base};
\end{axis}
\end{tikzpicture}
\end{column}
\begin{column}{0.38\textwidth}
\vspace{1mm}
\numero{1,87\vezes}{menor \ppl{} com a análise de sensibilidade, na mesma taxa de
poda: 56,08 contra 104,78.}
\vspace{3mm}
\numero{$\approx$\,19\%}{menos energia e tempo de inferência que o modelo base
retreinado, nas duas configurações ($\Delta E$ de 17,7\% e 19,1\%).}
\end{column}
\end{columns}
\vspace{1mm}
\nota{\textit{Perplexity} em escala logarítmica. Poda estruturada a 30\%, com e sem
análise de sensibilidade (\bet{1,0}); uma época no conjunto de treino do WikiText-2.}
\vfill
\mensagem{O dano medido sem retreinamento é parcialmente reversível, e o esquema de
poda com análise de sensibilidade continuou à frente ao mesmo custo.}
\end{frame}
